% !TeX root=../main.tex
\chapter{روش پیشنهادی بازسازی کور مشترک}
\label{chap:proposed-method}

\section{مقدمه}
این فصل قلب رساله است و معماری الگوریتم بازسازی کور مشترک کد توربو را از ایدهٔ اصلی تا تحلیل درستی شرح می‌دهد.

\section{ایده اصلی روش پیشنهادی}
ایدهٔ محوری، ترکیب قیود جبری RSC با آزمون آماری تحت نویز و جست‌وجوی درختی مقید است تا فضای جایگزشت‌ها و پلی‌نوم‌ها به‌صورت تدریجی هرس شود.

\section{تعریف فضای جست‌وجو}
\subsection{فضای پلی‌نوم‌ها}
\subsection{فضای حافظه}
\subsection{فضای درهم‌ریزها}
\subsection{محدودیت‌های طول فریم}
فضای خام درهم‌ریز از مرتبهٔ \(N!\) است و جست‌وجوی کامل عملاً ناممکن است؛ بنابراین کاهش فضا با قیود ضروری است.

\section{استخراج قیود جبری RSC}
\subsection{رابطه recursive}
\subsection{قیود parity}
\subsection{قیود polynomial realizability}
این قیود نامزدهای ناسازگار را در مراحل اولیه حذف می‌کنند.

\section{آزمون آماری قیود}
\subsection{احتمال برقراری constraint}
\subsection{احتمال نقض constraint}
\subsection{مدل Binomial}
\subsection{تعیین آستانه رد}
\subsection{کنترل احتمال پذیرش کاذب}
تحت نویز، برقراری دقیق قیود تضمین نمی‌شود؛ بنابراین آزمون آماری مبتنی بر مدل دوجمله‌ای برای پذیرش یا رد فرضیه‌ها به کار می‌رود.

\section{معماری جست‌وجوی درختی}
\subsection{تعریف Node}
\subsection{تعریف Depth}
\subsection{تعریف Hypothesis}
\subsection{Priority Queue}
\subsection{Best-first Search}
جست‌وجوی بهترین‌اول با صف اولویت، بسط فرضیه‌های امیدوارکننده‌تر را در اولویت قرار می‌دهد.

\section{فاز اول بازسازی}
\subsection{بازسازی \(U_1\)}
\subsection{بازسازی \(P_1\)}
\subsection{بازسازی \(q(D)\)}
\subsection{استخراج \(g_0(D)\)}
\subsection{استخراج \(g_1(D)\)}
\subsection{تعیین حافظه \(m\)}
\subsection{کنترل طول فریم}
فاز اول عمدتاً مؤلفه‌های مرتبط با کدگذار اول و طول سازگار را بازیابی می‌کند.

\section{فاز دوم بازسازی}
\subsection{بازسازی \(U_2\)}
\subsection{بازسازی \(P_2\)}
\subsection{قیود interleaver}
\subsection{بازیابی \(\pi_i\)}
\subsection{بازیابی \(\pi_o\)}
فاز دوم با اتکا به نتایج فاز اول، دنباله‌ها و درهم‌ریزهای داخلی و خارجی را بازسازی می‌کند.

\section{راهبردهای Pruning}
\subsection{Parity-consistency}
\subsection{Polynomial realizability}
\subsection{Interleaver consistency}
\subsection{Encoder-domain constraint}
\subsection{Heuristic pruning}
هرس چندلایه مانع انفجار شاخه‌ها می‌شود و باید احتمال حذف پاسخ واقعی را کنترل کند.

\section{الگوریتم نهایی}
در این بخش شبه‌کد کامل الگوریتم دوفازی همراه با ورودی/خروجی و معیار توقف ارائه می‌شود.

\begin{algorithm}[ht]
\caption{اسکلت الگوریتم بازسازی کور مشترک کد توربو}
\label{alg:joint-blind}
\begin{algorithmic}[1]
\REQUIRE مشاهدات \(R\)، محدودهٔ جست‌وجوی \(m\) و طول، آستانه‌های آماری
\ENSURE مجموعهٔ پاسخ‌های مجاز \(\widehat{\Theta}\)
\STATE مقداردهی صف اولویت با فرضیه‌های اولیهٔ فاز اول
\WHILE{صف خالی نیست}
	\STATE انتخاب بهترین فرضیه
	\STATE آزمون قیود جبری و آماری
	\IF{فرضیه رد شود}
		\STATE هرس و ادامه
	\ENDIF
	\STATE بسط فرزندان و به‌روزرسانی صف
\ENDWHILE
\STATE اجرای فاز دوم برای نامزدهای باقی‌مانده و تشکیل \(\widehat{\Theta}\)
\RETURN \(\widehat{\Theta}\)
\end{algorithmic}
\end{algorithm}

\section{تحلیل درستی الگوریتم}
\subsection{شرط لازم}
\subsection{شرط کافی}
\subsection{احتمال حذف جواب واقعی}
\subsection{احتمال پذیرش جواب کاذب}
\subsection{کامل بودن یا عدم کامل بودن جست‌وجو}
در این بخش، خواص نظری الگوریتم نسبت به مقالهٔ پایه گسترش می‌یابد تا برای سطح رسالهٔ دکتری کفایت استدلالی داشته باشد.
